\chapter{Planification prévisionnelle du projet (Gantt) et chemin critique}

\section{Planification par phases}

La planification s'étend sur six semaines de réalisation, avec une septième semaine consacrée exclusivement à la préparation et à la présentation du projet.

Les phases retenues sont les suivantes :

\begin{enumerate}
  \item \textbf{Analyse} (T1) : identification et validation des besoins, analyse des acteurs, rédaction des exigences et définition du périmètre.
  \item \textbf{Conception} (T2) : préparation des éléments de conception.
  \item \textbf{Développement backend} (T3 -- T4) : développement de la partie serveur (authentification, gestion des utilisateurs, annonces, catalogue, recherche/filtres, disponibilité).
  \item \textbf{Développement frontend} (T5) : développement des interfaces et intégration avec le backend (contact externe, modération).
  \item \textbf{Tests et Validation} (T6) : tests des fonctionnalités, corrections et validation.
  \item \textbf{Déploiement et Maintenance} (T7) : déploiement et actions de maintenance préliminaires.
  \item \textbf{Présentation} (T8) : préparation et présentation du projet (Semaine 7).
\end{enumerate}

\section{Tâches, durées, prédécesseurs et responsables}

Le tableau ci-dessous présente le détail des tâches.

\begin{table}[ht]
\centering
\renewcommand{\arraystretch}{1.25}
\small
\begin{tabularx}{\textwidth}{>{\bfseries}c X c c X}
\toprule
\textbf{ID} & \textbf{Tâche} & \textbf{Durée estimée} & \textbf{Prédécesseurs} & \textbf{Responsable principal} \\
\midrule
T1 & Analyse : besoins, acteurs, exigences et périmètre & 1 semaine & -- & Product Owner \\
T2 & Conception & 1 semaine & T1 & Architecte \\
T3 & Développement backend -- Partie 1 & 1 semaine & T2 & Développeur Backend \\
T4 & Développement backend -- Partie 2 & 1 semaine & T3 & Développeur Backend \\
T5 & Développement frontend & 1 semaine & T4 & Développeur Frontend / Équipe \\
T6 & Tests et Validation & 1 semaine & T5 & Toute l'équipe \\
T7 & Déploiement et Maintenance & 1 semaine & T6 & Tech Lead / Équipe \\
T8 & Présentation : préparation et présentation du projet & 1 semaine (S7) & T7 & Scrum Master / Toute l'équipe \\
\bottomrule
\end{tabularx}
\caption{Tâches, durées, prédécesseurs et responsables}
\end{table}

\section{Planification prévisionnelle (Diagramme de Gantt)}

La figure ci-dessous présente la planification prévisionnelle sur 7 semaines (S1 -- S7).

\begin{ganttchart}[
  hgrid,
  vgrid,
  x unit=1.6cm,
  y unit chart=0.55cm,
  bar/.append style={fill=blue!30},
  bar incomplete/.append style={fill=blue!15},
  title/.append style={fill=gray!10},
  title label font=\small\bfseries,
  bar label font=\small,
  group label font=\small\bfseries
]{1}{7}
  \gantttitle{S1}{1}
  \gantttitle{S2}{1}
  \gantttitle{S3}{1}
  \gantttitle{S4}{1}
  \gantttitle{S5}{1}
  \gantttitle{S6}{1}
  \gantttitle{S7}{1} \\

  \ganttbar{T1 - Analyse}{1}{1} \\
  \ganttbar{T2 - Conception}{2}{2} \\
  \ganttbar{T3 - Dév. backend Partie 1}{3}{3} \\
  \ganttbar{T4 - Dév. backend Partie 2}{4}{4} \\
  \ganttbar{T5 - Dév. frontend}{5}{5} \\
  \ganttbar{T6 - Tests et Validation}{6}{6} \\
  \ganttbar{T7 - Déploiement et Maintenance}{6}{6} \\
  \ganttbar{T8 - Présentation}{7}{7}
\end{ganttchart}
\caption{Planning prévisionnel (Gantt) sur 7 semaines (S1--S7)}
\label{gantt:planning}

\section{Chemin critique}

Le chemin critique déterminant la durée minimale du projet est le suivant :

\begin{center}
\textbf{T1 $\to$ T2 $\to$ T3 $\to$ T4 $\to$ T5 $\to$ T6 $\to$ T7 $\to$ T8}
\end{center}

\textbf{Interprétation :} tout retard sur une tâche appartenant à ce chemin critique entraînerait, sans réorganisation du travail, un retard équivalent sur la livraison finale.
